\chapter{Status and Anchor Summary}
\label{chap:reader-status-anchor-summary}

This chapter keeps only the minimum global navigation state needed for the
reader build: the six status tags and the top-level result-anchor families.

\section{Status tags}

\begin{longtable}{L{0.25\textwidth}L{0.67\textwidth}}
\toprule
Status tag & Reader-build meaning \\
\midrule
\endhead
\StatusExact & Direct consequence of the declared parent theory, exact
definitions, or exact algebra with no extra closure step. \\
\addlinespace
\StatusExactClosure & Exact once a named closure, branch family, or search
class has been declared. \\
\addlinespace
\StatusReduced & Valid only after a controlled reduction, asymptotic regime, or
selected branch restriction. \\
\addlinespace
\StatusEffective & Organized by a useful model choice or closure that is not yet
uniquely forced by the full PDE. \\
\addlinespace
\StatusNumerical & Defined by exact equations but currently located or checked
numerically. \\
\addlinespace
\StatusOpen & Still depends on actual branch realization or another missing
theorem gate. \\
\bottomrule
\end{longtable}

\section{Top-level result-anchor families}

\begin{longtable}{L{0.13\textwidth}L{0.14\textwidth}L{0.62\textwidth}}
\toprule
Anchor & Stages & Main role in the derivation program \\
\midrule
\endhead
\resultanchor{MTDC-T1} & Parent + 1 & Parent-theory import, projection and
reduction language, mixed-sector firewall, and first moving-throat target
statement. \\
\addlinespace
\resultanchor{MTDC-T2} & 1--2 & Moving-surface lift, recovery of \((a,L)\), and
distributed wall geometry closure. \\
\addlinespace
\resultanchor{MTDC-T3} & 3 & Stable BdG reduction, Schur-complement support
kernels, and first grouped \(P_2\) conservative diagnostics. \\
\addlinespace
\resultanchor{MTDC-T4} & 4--5 & Maxwell/mixed reduction, outgoing bridge, and
the invariant outgoing-normalization product. \\
\addlinespace
\resultanchor{MTDC-T5} & 6--19 & Full grouped bundle, overlap and isotropy
machinery, selected-mode normalization language, and support-feasibility front
end. \\
\addlinespace
\resultanchor{MTDC-T6} & 20--34 & Continuum placement, split-\(U\) branch,
rank-2 support completion, and coherent tracking front end. \\
\addlinespace
\resultanchor{MTDC-T7} & 35--73 & Gain thresholds, Family-1 branch windows, and
the first reduced Family-1 verdict chain. \\
\addlinespace
\resultanchor{MTDC-T8} & 74--146 & Geometry-lane firewall, retarded closure,
outlet/core/mouth compensation, and renormalized normalization residuals. \\
\addlinespace
\resultanchor{MTDC-T9} & 147--183 & Weak-axisymmetric transport, quotient
coordinates, orbit closure, and the four-scalar verdict packet. \\
\addlinespace
\resultanchor{MTDC-T10} & 184--201 & Realization compiler, graph-error form,
finite search sieve, and local mixed-ray closure. \\
\addlinespace
\resultanchor{MTDC-T11} & 202--225 & Same-charge audit, rigid-mouth chart,
selected twin-support placement, and strict actual-branch packet. \\
\addlinespace
\resultanchor{MTDC-T12} & 226--236 & Relaxed open-system branch, recovery map,
barrier/export companions, and material-threshold layer. \\
\bottomrule
\end{longtable}

\section{Status discipline}

The reader build uses the same weakest-link rule as the archive build: a result
may only be cited at the strongest level supported by its weakest essential
input. When in doubt, the archive build and the verification trackers control
the stronger interpretation.

\section{Verification routing}

The stage number is the reader-facing audit handle.  This build intentionally
omits the full verification filename block from each stage page and uses a
shared naming convention instead.

All paths below are relative to \nolinkurl{research/pde_ledger/}.

\begin{itemize}[leftmargin=2.2em]
  \item SymPy symbolic audits follow
  \nolinkurl{moving_throat_pde_stageNNN_*_sympy_audit.py}.
  \item Mathematica symbolic audits follow
  \nolinkurl{moving_throat_pde_stageNNN_*_mathematica_audit.wl}.
  \item Numerical stress harnesses usually follow
  \nolinkurl{stageNNN*_stress.py} and \nolinkurl{stageNNN*_stress.wl}, with
  occasional shared-range names such as \nolinkurl{stage003_004_*}.
  \item Exact path inventory and trust status live in the repo verification
  trackers rather than in each printed stage.
\end{itemize}
